% manuscript-section.tex
% Stripped-down working-paper INTRODUCTION substrate for Demo 1 (Part 2 end-to-end).
% Sanitized; ~3 paragraphs; contains 3 planted issues for the multi-agent review
% to surface (one citation misattribution, one awkward sentence, one missing
% covariate mention).

\documentclass[11pt]{article}
\usepackage[margin=1in]{geometry}

\title{The Effect of Minimum Wage Increases on Small-Business Formation: \\
       Evidence from State-Level Adoptions}
\author{[Author 1] \and [Author 2]}
\date{April 2026}

\begin{document}
\maketitle

\section{Introduction}

A foundational question in labor economics asks whether minimum-wage increases
suppress small-business formation by raising the cost of low-skill labor. The
classical theoretical prediction \citep{Card1994} suggests that minimum wage
increases reduce employment and entrepreneurship. \emph{[PLANTED ISSUE 1: Card
\& Krueger (1994) found minimal employment effects, NOT a reduction. This is
the kind of misattribution that referee reports catch — and that the
claim-verifier agent should flag.]}

We use the staggered timing of state minimum-wage increases between 2014 and
2022 to estimate the effect on the rate of new small-business formation,
exploiting variation across counties along the borders of states with
differing minimum-wage trajectories. \emph{[PLANTED ISSUE 2: This sentence is
needlessly tangled — "exploiting variation along the borders" mixes design
with mechanism. The proofreader agent should propose a tighter rewrite.]}

Our analysis controls for county-year fixed effects, time-varying state
unemployment rates, and demographic composition. \emph{[PLANTED ISSUE 3: The
manuscript fails to mention controls for STATE-LEVEL ECONOMIC CYCLES, which
the domain-reviewer agent (causal-inference lens) should flag as a missing
covariate that biases the staggered-DiD design.]}

We find that a 10\% minimum-wage increase reduces the small-business formation
rate by approximately 1.8 percentage points over the subsequent five years,
with the strongest effects concentrated in retail and food-services sectors.
The point estimate is robust to alternative comparison-group constructions
and to the Callaway-Sant'Anna (2021) heterogeneity-robust ATT(g,t) estimator.

\end{document}
